<!-- Generated by pkgdown: do not edit by hand -->
<!DOCTYPE html>
<html>
  <head>
  <meta charset="utf-8">
<meta http-equiv="X-UA-Compatible" content="IE=edge">
<meta name="viewport" content="width=device-width, initial-scale=1.0">

<title>Data Analysis with \texttt{augmentedRCBD} • augmentedRCBD</title>

<!-- jquery -->
<script src="https://code.jquery.com/jquery-3.1.0.min.js" integrity="sha384-nrOSfDHtoPMzJHjVTdCopGqIqeYETSXhZDFyniQ8ZHcVy08QesyHcnOUpMpqnmWq" crossorigin="anonymous"></script>
<!-- Bootstrap -->
<link href="https://maxcdn.bootstrapcdn.com/bootswatch/3.3.7/united/bootstrap.min.css" rel="stylesheet" crossorigin="anonymous">

<script src="https://maxcdn.bootstrapcdn.com/bootstrap/3.3.7/js/bootstrap.min.js" integrity="sha384-Tc5IQib027qvyjSMfHjOMaLkfuWVxZxUPnCJA7l2mCWNIpG9mGCD8wGNIcPD7Txa" crossorigin="anonymous"></script>

<!-- Font Awesome icons -->
<link href="https://maxcdn.bootstrapcdn.com/font-awesome/4.6.3/css/font-awesome.min.css" rel="stylesheet" integrity="sha384-T8Gy5hrqNKT+hzMclPo118YTQO6cYprQmhrYwIiQ/3axmI1hQomh7Ud2hPOy8SP1" crossorigin="anonymous">

<!-- clipboard.js -->
<script src="https://cdnjs.cloudflare.com/ajax/libs/clipboard.js/1.7.1/clipboard.min.js" integrity="sha384-cV+rhyOuRHc9Ub/91rihWcGmMmCXDeksTtCihMupQHSsi8GIIRDG0ThDc3HGQFJ3" crossorigin="anonymous"></script>

<!-- sticky kit -->
<script src="https://cdnjs.cloudflare.com/ajax/libs/sticky-kit/1.1.3/sticky-kit.min.js" integrity="sha256-c4Rlo1ZozqTPE2RLuvbusY3+SU1pQaJC0TjuhygMipw=" crossorigin="anonymous"></script>

<!-- pkgdown -->
<link href="../pkgdown.css" rel="stylesheet">
<script src="../pkgdown.js"></script>


<!-- docsearch -->
<script src="../docsearch.js"></script>
<link rel="stylesheet" href="https://cdn.jsdelivr.net/npm/docsearch.js@2/dist/cdn/docsearch.min.css" />
<link href="../docsearch.css" rel="stylesheet">
<script src="https://cdnjs.cloudflare.com/ajax/libs/mark.js/8.11.1/jquery.mark.min.js" integrity="sha256-4HLtjeVgH0eIB3aZ9mLYF6E8oU5chNdjU6p6rrXpl9U=" crossorigin="anonymous"></script>


  <link href="../extra.css" rel="stylesheet">
  
<meta property="og:title" content="Data Analysis with \texttt{augmentedRCBD}" />

<meta property="og:description" content="" />
<meta name="twitter:card" content="summary" />



<!-- mathjax -->
<script src='https://mathjax.rstudio.com/latest/MathJax.js?config=TeX-AMS-MML_HTMLorMML'></script>

<!--[if lt IE 9]>
<script src="https://oss.maxcdn.com/html5shiv/3.7.3/html5shiv.min.js"></script>
<script src="https://oss.maxcdn.com/respond/1.4.2/respond.min.js"></script>
<![endif]-->


<!-- Global site tag (gtag.js) - Google Analytics -->
<script async src="https://www.googletagmanager.com/gtag/js?id=UA-116716530-2"></script>
<script>
  window.dataLayer = window.dataLayer || [];
  function gtag(){dataLayer.push(arguments);}
  gtag('js', new Date());

  gtag('config', 'UA-116716530-2');
</script>

  </head>

  <body>
    <div class="container template-article">
      <header>
      <div class="navbar navbar-default navbar-fixed-top" role="navigation">
  <div class="container">
    <div class="navbar-header">
      <button type="button" class="navbar-toggle collapsed" data-toggle="collapse" data-target="#navbar">
        <span class="icon-bar"></span>
        <span class="icon-bar"></span>
        <span class="icon-bar"></span>
      </button>
      <span class="navbar-brand">
        <a class="navbar-link" href="../index.html">augmentedRCBD</a>
        <span class="label label-danger" data-toggle="tooltip" data-placement="bottom" title="Unreleased package">0.1.0.9000</span>
      </span>
    </div>

    <div id="navbar" class="navbar-collapse collapse">
      <ul class="nav navbar-nav">
        <li>
  <a href="../index.html">
    <span class="fa fa-home fa-lg"></span>
     
    Home
  </a>
</li>
<li>
  <a href="../reference/index.html">
    <span class="fa fa-code"></span>
     
    Reference
  </a>
</li>
<li>
  <a href="../articles/Data_Analysis_with_augmentedRCBD.html">
    <span class="fa fa-file-o"></span>
     
    Tutorial
  </a>
</li>
<li>
  <a href="../news/index.html">
    <span class="fa fa-newspaper-o"></span>
     
    News
  </a>
</li>
      </ul>
      
      <ul class="nav navbar-nav navbar-right">
        <li>
  <a href="https://github.com/aravind-j/augmentedRCBD">
    <span class="fa fa-github-alt"></span>
     
    Github
  </a>
</li>
      </ul>
      
      <form class="navbar-form navbar-right" role="search">
        <div class="form-group">
          <input type="search" class="form-control" name="search-input" id="search-input" placeholder="Search..." aria-label="Search for..." autocomplete="off">
        </div>
      </form>
      
    </div><!--/.nav-collapse -->
  </div><!--/.container -->
</div><!--/.navbar -->

      
      </header>

\usepackage{fancyhdr}
\usepackage{wrapfig}
\usepackage{float}
\pagestyle{fancy}
\fancyhead[LE,RO]{\slshape \rightmark}
\fancyhead[LO,RE]{Data Analysis with \texttt{augmentedRCBD}}
\fancyfoot[C]{\thepage}
\usepackage{hyperref}
\hypersetup{colorlinks=true}
\hypersetup{linktoc=all}
\hypersetup{linkcolor=blue}
\usepackage{pdflscape}
\usepackage{booktabs}
\newcommand{\blandscape}{\begin{landscape}}
\newcommand{\elandscape}{\end{landscape}}
\renewcommand\thesection{\arabic{section}}

<div class="row">
  <div class="col-md-9 contents">
    <div class="page-header toc-ignore">
      <h1>Data Analysis with \texttt{augmentedRCBD}</h1>
                        <h4 class="author">Aravind, J.\textsuperscript{1}, Mukesh Sankar, S.\textsuperscript{2},
Wankhede, D. P.\textsuperscript{3}, and Kaur, V.\textsuperscript{4}</h4>
            
            <h4 class="date">2019-01-17</h4>
      
      <small class="dont-index">Source: <a href='https://github.com/aravind-j/augmentedRCBD
    https://CRAN.R-project.org/package=augmentedRCBD
    https://aravind-j.github.io/augmentedRCBD
    https://doi.org/10.5281/zenodo.1310011/blob/master/vignettes/Data_Analysis_with_augmentedRCBD.Rmd'><code>vignettes/Data_Analysis_with_augmentedRCBD.Rmd</code></a></small>
      <div class="hidden name"><code>Data_Analysis_with_augmentedRCBD.Rmd</code></div>

    </div>

    
    
\begin{center}
1. Division of Germplasm Conservation, ICAR-National Bureau of Plant Genetic Resources, New Delhi.

2. Division of Genetics, ICAR-Indian Agricultural Research Institute, New Delhi.

3. Division of Genomic Resources, ICAR-National Bureau of Plant Genetic Resources, New Delhi.

4. Division of Germplasm Evaluation, ICAR-National Bureau of Plant Genetic Resources, New Delhi.

\end{center}

\begin{center}
\vspace{6pt}
\hrule
\end{center}

\tableofcontents

\begin{wrapfigure}{r}{0.35\textwidth}
  \vspace{-10pt}
  \begin{center}
    \includegraphics[width=0.33\textwidth]{C:/Users/lenovo/Documents/R/win-library/3.6/augmentedRCBD/extdata/augmentedRCBD.png}
  \end{center}
    \vspace{-10pt}
\end{wrapfigure}

\hypertarget{overview}{%
\section{Overview}\label{overview}}

The software \texttt{augmentedRCBD} is built on the
\href{https://en.wikipedia.org/wiki/R_(programming_language)}{\texttt{R}
statistical programming language} as an add-on (or `package' in the
\texttt{R} \emph{lingua franca}). It performs the analysis of data
generated from experiments in augmented randomised complete block design
according to Federer, W.T.
(\protect\hyperlink{ref-federer_augmented_1956}{1956},
\protect\hyperlink{ref-federer_augmented_1961}{1961}). It also computes
analysis of variance, adjusted means, descriptive statistics, genetic
variability statistics etc. and includes options for data visualization
and report generation.

This tutorial aims to educate the users in utilising this package for
performing such analysis. Utilising \texttt{augmentedRCBD} for data
analysis requires a basic knowledge of \texttt{R} programming language.
However, as many of the intended end-users may not be familiar with
\texttt{R}, \protect\hyperlink{rsoft}{sections 2 to 4} give a `gentle'
introduction to \texttt{R}, especially those aspects which are necessary
to get \texttt{augmentedRCBD} up and running for performing data
analysis in a Windows environment. Users already familiar with
\texttt{R} can feel free to skip to \protect\hyperlink{install}{section
5}.

\begin{wrapfigure}{r}{0.35\textwidth}
  \vspace{-10pt}
  \begin{center}
    \includegraphics[width=0.20\textwidth]{rlogo.png}
  \end{center}
    \vspace{-10pt}
\end{wrapfigure}

\hypertarget{rsoft}{%
\section{\texorpdfstring{\texttt{R} software}{R software}}\label{rsoft}}

It is a free software environment for statistical computing and
graphics. It is free and open source, platform independent (works on
Linux, Windows or MacOS), very flexible, comprehensive with robust
interfaces for all the popular programming languates as well as
databases. It is strengthened by its diverse library of add-on packages
extending its ability as well as the incredible community support. It is
one of the most popular tools being used in academia today (Tippmann,
\protect\hyperlink{ref-tippmann_programming_2015}{2015}).

\hypertarget{getting-started}{%
\section{Getting Started}\label{getting-started}}

This section details the steps required to set up the \texttt{R}
programming environment under a third-party interface called
\texttt{RStudio} in Windows.

\hypertarget{installing-r}{%
\subsection{\texorpdfstring{Installing
\texttt{R}}{Installing R}}\label{installing-r}}

Download and install \texttt{R} for Windows from
\url{http://cran.r-project.org/bin/windows/base/}.

\includegraphics{rbase.png}

\begin{center}
                \textbf{Fig. 1}: The \texttt{R} download location.
                \end{center}

\hypertarget{installing-rstudio}{%
\subsection{\texorpdfstring{Installing
\texttt{RStudio}}{Installing RStudio}}\label{installing-rstudio}}

The basic
\href{https://en.wikipedia.org/wiki/Command_line_interface}{command line
interface} in native \texttt{R} is rather limiting. There are several
interfaces which enhance it's functionality and ease of use,
\href{https://www.rstudio.com/}{\texttt{RStudio}} being one of the most
popular among \texttt{R} programmers.

Download and install \texttt{RStudio} for Windows from
\url{https://www.rstudio.com/products/rstudio/download/\#download}

\includegraphics{rstudio.png}

\begin{center}
                \textbf{Fig. 2}: The \texttt{RStudio} download location.
                \end{center}

\hypertarget{the-rstudio-interface}{%
\subsection{\texorpdfstring{The \texttt{RStudio}
Interface}{The RStudio Interface}}\label{the-rstudio-interface}}

On opening \texttt{RStudio}, the default interface with four
panes/windows is visible as follows. Few panes have different tabs.

\includegraphics{rstudio panes.png}

\begin{center}
                \textbf{Fig. 3}: The default \texttt{RStudio} interface with the four panes.
                \end{center}

\hypertarget{console}{%
\subsubsection{Console}\label{console}}

This is where the action happens. Here any authentic \texttt{R} code
typed after the `\texttt{\textgreater{}}' prompt will be executed after
pressing `Enter' to generate the output.

For example, type \texttt{1+1} in the console and press `Enter'.

\begin{Shaded}
\begin{Highlighting}[]
\DecValTok{1}\OperatorTok{+}\DecValTok{1}
\end{Highlighting}
\end{Shaded}

\begin{verbatim}
[1] 2
\end{verbatim}

\hypertarget{source}{%
\subsubsection{Source}\label{source}}

This is where \texttt{R} Scripts (collection of code) can be created and
edited. \texttt{R} scripts are text files with a \texttt{.R} extension.
\texttt{R} Code for analysis can be typed and saved in such \texttt{R}
scripts. New scripts can be opened by clicking `File\textbar{}New File'
and selecting `R Script'. Code can be selected from \texttt{R} Scripts
and sent to console for evaluation by clicking `Run' on the `Source'
pane or by pressing `Ctrl + Enter'.

\hypertarget{environmenthistoryconnections}{%
\subsubsection{Environment\textbar{}History\textbar{}Connections}\label{environmenthistoryconnections}}

The `Environment' tab shows the list of all the `objects' (see
\protect\hyperlink{ObjFun}{section 4.3}) defined in the current
\texttt{R} session. It has also some buttons up top to open, save and
clear the environment as well as few options for import of data under
\texttt{Import\ Dataset}.

The `History' tab shows a history of all the code that was previously
evaluated. This is useful, if you want to go back to some code.

The `Connections' tab helps to establish and manage connections with
different databases and data sources.

\hypertarget{filesplotspackageshelpviewer}{%
\subsubsection{Files\textbar{}Plots\textbar{}Packages\textbar{}Help\textbar{}Viewer}\label{filesplotspackageshelpviewer}}

The `Files' tab shows a sleek file browser to access the file directory
in the computer with options to manage the working directory (see
\protect\hyperlink{wdir}{section 4.1}) under the More button.

The `Plots' tab shows all the plots generated in \texttt{R} with buttons
to delete unnecessary ones and export useful ones as a pdf file or as an
image file.

The `Packages' tab shows a list of all the \texttt{R} add-on packages
installed. The check box on the left shows whether they are loaded or
not. There are also buttons to install and update \texttt{R} packages.

The `Viewer' tab shows any web content output generated by an \texttt{R}
code.

\hypertarget{some-basics}{%
\section{Some Basics}\label{some-basics}}

This section describes some basics to enable the users to have a working
knowledge in \texttt{R} in order to use \texttt{augmentedRCBD}.

\hypertarget{wdir}{%
\subsection{Working Directory}\label{wdir}}

It is a file path to a folder on the computer which is recognised by
\texttt{R} as the default location to read files from or write files to.
The code \texttt{getwd()} shows the current working directory, while
\texttt{setwd()} can be used to change the existing working directory.

\begin{Shaded}
\begin{Highlighting}[]
\CommentTok{# Print current working directory}
\KeywordTok{getwd}\NormalTok{()}
\end{Highlighting}
\end{Shaded}

\begin{verbatim}
[1] "C:/Users/Computer/Documents"
\end{verbatim}

\begin{Shaded}
\begin{Highlighting}[]
\CommentTok{# Set new working directory}
\KeywordTok{setwd}\NormalTok{(}\StringTok{"C:/Data Analysis/"}\NormalTok{)}
\KeywordTok{getwd}\NormalTok{()}
\end{Highlighting}
\end{Shaded}

\begin{verbatim}
[1] "C:/Data Analysis/"
\end{verbatim}

One key detail is that file paths in \texttt{R} uses forward slashes
(\texttt{/}) as in MacOS or Linux, unlike backward slashes
(\texttt{\textbackslash{}}) in Windows. This needs to be considered
while copying paths from default Windows file explorer.

\hypertarget{expression-and-assignment}{%
\subsection{Expression and Assignment}\label{expression-and-assignment}}

Expressions are instructions in the form of code to be entered after the
\texttt{\textgreater{}} prompt in the console. Expressions can be a
constant, an arithmetic or a condition. A more advanced and most useful
expression is a function call (see \protect\hyperlink{ObjFun}{section
4.3}).

\begin{Shaded}
\begin{Highlighting}[]
\CommentTok{# Constant}
\DecValTok{123}
\end{Highlighting}
\end{Shaded}

\begin{verbatim}
[1] 123
\end{verbatim}

\begin{Shaded}
\begin{Highlighting}[]
\CommentTok{# Arithmetic (add two numbers)}
\DecValTok{1} \OperatorTok{+}\StringTok{ }\DecValTok{2}
\end{Highlighting}
\end{Shaded}

\begin{verbatim}
[1] 3
\end{verbatim}

\begin{Shaded}
\begin{Highlighting}[]
\CommentTok{# Condition}
\DecValTok{34} \OperatorTok{>}\StringTok{ }\DecValTok{25}
\end{Highlighting}
\end{Shaded}

\begin{verbatim}
[1] TRUE
\end{verbatim}

\begin{Shaded}
\begin{Highlighting}[]
\DecValTok{1} \OperatorTok{==}\StringTok{ }\DecValTok{2}
\end{Highlighting}
\end{Shaded}

\begin{verbatim}
[1] FALSE
\end{verbatim}

\begin{Shaded}
\begin{Highlighting}[]
\CommentTok{# Function call (mean of a series of numbers)}
\KeywordTok{mean}\NormalTok{(}\KeywordTok{c}\NormalTok{(}\DecValTok{25}\NormalTok{,}\DecValTok{56}\NormalTok{,}\DecValTok{89}\NormalTok{,}\DecValTok{35}\NormalTok{))}
\end{Highlighting}
\end{Shaded}

\begin{verbatim}
[1] 51.25
\end{verbatim}

Information from an expression can be stored as an `object' (see
\protect\hyperlink{ObjFun}{section 4.3}) by assigning a name using the
operator `\texttt{\textless{}-}'.

\begin{Shaded}
\begin{Highlighting}[]
\CommentTok{# Assign the result of the expression 1 + 2 to an object 'a'}
\NormalTok{a <-}\StringTok{ }\DecValTok{1} \OperatorTok{+}\StringTok{ }\DecValTok{2}
\NormalTok{a}
\end{Highlighting}
\end{Shaded}

\begin{verbatim}
[1] 3
\end{verbatim}

It is recommended to add comments to explain the code by using the
`\texttt{\#}' sign. Any code after the `\texttt{\#}' sign will be
ignored by \texttt{R}.

\hypertarget{ObjFun}{%
\subsection{Objects and Functions}\label{ObjFun}}

\texttt{R} is an object-oriented programming language (OOP). Any kind or
construct created in \texttt{R} is an `object'. Each object has a
`class' (shown using the \texttt{class()} function) and different
`attributes' which defines what operations can be done on that object.
There are different types of data structure objects in \texttt{R} such
as vectors, matrices, factors, data frames, and lists. A `function' is
also an object, which defines a procedure or a sequence of expressions.

\hypertarget{vector}{%
\subsubsection{Vector}\label{vector}}

A vector is a collection of elements of a single type (or `mode'). The
common vector modes are `numeric', `integer', `character' and `logical'.
The \texttt{c()} function is used to create vectors. The functions
\texttt{class()}, \texttt{str()} and \texttt{length()} show the
attributes of vectors.

Vector modes `numeric' stores real numbers, while `integer' stores
integers, which can be enforced by suffixing elements with `\texttt{L}'.

\begin{Shaded}
\begin{Highlighting}[]
\CommentTok{# A numeric vector}
\NormalTok{a <-}\StringTok{ }\KeywordTok{c}\NormalTok{(}\DecValTok{1}\NormalTok{, }\DecValTok{2}\NormalTok{, }\FloatTok{3.3}\NormalTok{)}
\KeywordTok{class}\NormalTok{(a)}
\end{Highlighting}
\end{Shaded}

\begin{verbatim}
[1] "numeric"
\end{verbatim}

\begin{Shaded}
\begin{Highlighting}[]
\KeywordTok{str}\NormalTok{(a)}
\end{Highlighting}
\end{Shaded}

\begin{verbatim}
 num [1:3] 1 2 3.3
\end{verbatim}

\begin{Shaded}
\begin{Highlighting}[]
\KeywordTok{length}\NormalTok{(a)}
\end{Highlighting}
\end{Shaded}

\begin{verbatim}
[1] 3
\end{verbatim}

\begin{Shaded}
\begin{Highlighting}[]
\CommentTok{# An integer vector}
\NormalTok{b <-}\StringTok{ }\KeywordTok{c}\NormalTok{(1L, 2L, 3L)}
\KeywordTok{class}\NormalTok{(b)}
\end{Highlighting}
\end{Shaded}

\begin{verbatim}
[1] "integer"
\end{verbatim}

\begin{Shaded}
\begin{Highlighting}[]
\KeywordTok{str}\NormalTok{(b)}
\end{Highlighting}
\end{Shaded}

\begin{verbatim}
 int [1:3] 1 2 3
\end{verbatim}

\begin{Shaded}
\begin{Highlighting}[]
\KeywordTok{length}\NormalTok{(b)}
\end{Highlighting}
\end{Shaded}

\begin{verbatim}
[1] 3
\end{verbatim}

The vector mode `character' store text.

\begin{Shaded}
\begin{Highlighting}[]
\CommentTok{# A character vector}
\NormalTok{c <-}\StringTok{ }\KeywordTok{c}\NormalTok{(}\StringTok{"one"}\NormalTok{,}\StringTok{"two"}\NormalTok{,}\StringTok{"three"}\NormalTok{)}
\KeywordTok{class}\NormalTok{(c)}
\end{Highlighting}
\end{Shaded}

\begin{verbatim}
[1] "character"
\end{verbatim}

\begin{Shaded}
\begin{Highlighting}[]
\KeywordTok{str}\NormalTok{(c)}
\end{Highlighting}
\end{Shaded}

\begin{verbatim}
 chr [1:3] "one" "two" "three"
\end{verbatim}

\begin{Shaded}
\begin{Highlighting}[]
\KeywordTok{length}\NormalTok{(c)}
\end{Highlighting}
\end{Shaded}

\begin{verbatim}
[1] 3
\end{verbatim}

The vector mode `logical' stores `\texttt{TRUE}' OR `\texttt{FALSE}'
logical data.

\begin{Shaded}
\begin{Highlighting}[]
\CommentTok{#logical vector}
\NormalTok{d <-}\StringTok{ }\KeywordTok{c}\NormalTok{(}\OtherTok{TRUE}\NormalTok{,}\OtherTok{TRUE}\NormalTok{,}\OtherTok{TRUE}\NormalTok{,}\OtherTok{FALSE}\NormalTok{,}\OtherTok{TRUE}\NormalTok{,}\OtherTok{FALSE}\NormalTok{)}
\KeywordTok{class}\NormalTok{(d)}
\end{Highlighting}
\end{Shaded}

\begin{verbatim}
[1] "logical"
\end{verbatim}

\begin{Shaded}
\begin{Highlighting}[]
\KeywordTok{str}\NormalTok{(d)}
\end{Highlighting}
\end{Shaded}

\begin{verbatim}
 logi [1:6] TRUE TRUE TRUE FALSE TRUE FALSE
\end{verbatim}

\begin{Shaded}
\begin{Highlighting}[]
\KeywordTok{length}\NormalTok{(d)}
\end{Highlighting}
\end{Shaded}

\begin{verbatim}
[1] 6
\end{verbatim}

\hypertarget{factor}{%
\subsubsection{Factor}\label{factor}}

A `factor' in \texttt{R} stores data from categorical data in variables
as different levels.

\begin{Shaded}
\begin{Highlighting}[]
\NormalTok{catg <-}\StringTok{ }\KeywordTok{c}\NormalTok{(}\StringTok{"male"}\NormalTok{,}\StringTok{"female"}\NormalTok{,}\StringTok{"female"}\NormalTok{,}\StringTok{"male"}\NormalTok{,}\StringTok{"male"}\NormalTok{)}
\NormalTok{catg}
\end{Highlighting}
\end{Shaded}

\begin{verbatim}
[1] "male"   "female" "female" "male"   "male"  
\end{verbatim}

\begin{Shaded}
\begin{Highlighting}[]
\KeywordTok{is.factor}\NormalTok{(catg)}
\end{Highlighting}
\end{Shaded}

\begin{verbatim}
[1] FALSE
\end{verbatim}

\begin{Shaded}
\begin{Highlighting}[]
\CommentTok{# Apply the factor function}
\NormalTok{factor_catg <-}\StringTok{ }\KeywordTok{factor}\NormalTok{(catg)}

\NormalTok{factor_catg}
\end{Highlighting}
\end{Shaded}

\begin{verbatim}
[1] male   female female male   male  
Levels: female male
\end{verbatim}

\begin{Shaded}
\begin{Highlighting}[]
\KeywordTok{is.factor}\NormalTok{(factor_catg)}
\end{Highlighting}
\end{Shaded}

\begin{verbatim}
[1] TRUE
\end{verbatim}

\begin{Shaded}
\begin{Highlighting}[]
\KeywordTok{class}\NormalTok{(factor_catg)}
\end{Highlighting}
\end{Shaded}

\begin{verbatim}
[1] "factor"
\end{verbatim}

\begin{Shaded}
\begin{Highlighting}[]
\KeywordTok{str}\NormalTok{(factor_catg)}
\end{Highlighting}
\end{Shaded}

\begin{verbatim}
 Factor w/ 2 levels "female","male": 2 1 1 2 2
\end{verbatim}

A character, numeric or integer vector can be transformed to a factor by
using the \texttt{as.factor()} function.

\begin{Shaded}
\begin{Highlighting}[]
\CommentTok{# Conversion of numeric to factor}
\NormalTok{a <-}\StringTok{ }\KeywordTok{c}\NormalTok{(}\DecValTok{1}\NormalTok{, }\DecValTok{2}\NormalTok{, }\FloatTok{3.3}\NormalTok{)}
\KeywordTok{class}\NormalTok{(a)}
\end{Highlighting}
\end{Shaded}

\begin{verbatim}
[1] "numeric"
\end{verbatim}

\begin{Shaded}
\begin{Highlighting}[]
\KeywordTok{str}\NormalTok{(a)}
\end{Highlighting}
\end{Shaded}

\begin{verbatim}
 num [1:3] 1 2 3.3
\end{verbatim}

\begin{Shaded}
\begin{Highlighting}[]
\NormalTok{fac_a <-}\StringTok{ }\KeywordTok{as.factor}\NormalTok{(a)}
\KeywordTok{class}\NormalTok{(fac_a)}
\end{Highlighting}
\end{Shaded}

\begin{verbatim}
[1] "factor"
\end{verbatim}

\begin{Shaded}
\begin{Highlighting}[]
\KeywordTok{str}\NormalTok{(fac_a)}
\end{Highlighting}
\end{Shaded}

\begin{verbatim}
 Factor w/ 3 levels "1","2","3.3": 1 2 3
\end{verbatim}

\begin{Shaded}
\begin{Highlighting}[]
\CommentTok{# Conversion of integer to factor}
\NormalTok{b <-}\StringTok{ }\KeywordTok{c}\NormalTok{(1L, 2L, 3L)}
\KeywordTok{class}\NormalTok{(b)}
\end{Highlighting}
\end{Shaded}

\begin{verbatim}
[1] "integer"
\end{verbatim}

\begin{Shaded}
\begin{Highlighting}[]
\KeywordTok{str}\NormalTok{(b)}
\end{Highlighting}
\end{Shaded}

\begin{verbatim}
 int [1:3] 1 2 3
\end{verbatim}

\begin{Shaded}
\begin{Highlighting}[]
\NormalTok{fac_b <-}\StringTok{ }\KeywordTok{as.factor}\NormalTok{(b)}
\KeywordTok{class}\NormalTok{(fac_b)}
\end{Highlighting}
\end{Shaded}

\begin{verbatim}
[1] "factor"
\end{verbatim}

\begin{Shaded}
\begin{Highlighting}[]
\KeywordTok{str}\NormalTok{(fac_b)}
\end{Highlighting}
\end{Shaded}

\begin{verbatim}
 Factor w/ 3 levels "1","2","3": 1 2 3
\end{verbatim}

\begin{Shaded}
\begin{Highlighting}[]
\CommentTok{# Conversion of character to factor}
\NormalTok{c <-}\StringTok{ }\KeywordTok{c}\NormalTok{(}\StringTok{"one"}\NormalTok{,}\StringTok{"two"}\NormalTok{,}\StringTok{"three"}\NormalTok{)}
\KeywordTok{class}\NormalTok{(c)}
\end{Highlighting}
\end{Shaded}

\begin{verbatim}
[1] "character"
\end{verbatim}

\begin{Shaded}
\begin{Highlighting}[]
\KeywordTok{str}\NormalTok{(c)}
\end{Highlighting}
\end{Shaded}

\begin{verbatim}
 chr [1:3] "one" "two" "three"
\end{verbatim}

\begin{Shaded}
\begin{Highlighting}[]
\NormalTok{fac_c <-}\StringTok{ }\KeywordTok{as.factor}\NormalTok{(c)}
\KeywordTok{class}\NormalTok{(fac_c)}
\end{Highlighting}
\end{Shaded}

\begin{verbatim}
[1] "factor"
\end{verbatim}

\begin{Shaded}
\begin{Highlighting}[]
\KeywordTok{str}\NormalTok{(fac_c)}
\end{Highlighting}
\end{Shaded}

\begin{verbatim}
 Factor w/ 3 levels "one","three",..: 1 3 2
\end{verbatim}

\hypertarget{matrix}{%
\subsubsection{Matrix}\label{matrix}}

A `matrix' in \texttt{R} is a vector with the attributes `\texttt{nrow}'
and `\texttt{ncol}'.

\begin{Shaded}
\begin{Highlighting}[]
\CommentTok{# Generate 5 * 4 numeric matrix}
\NormalTok{m <-}\StringTok{ }\KeywordTok{matrix}\NormalTok{(}\DecValTok{1}\OperatorTok{:}\DecValTok{20}\NormalTok{, }\DataTypeTok{nrow =} \DecValTok{5}\NormalTok{, }\DataTypeTok{ncol =} \DecValTok{4}\NormalTok{)}
\NormalTok{m}
\end{Highlighting}
\end{Shaded}

\begin{verbatim}
     [,1] [,2] [,3] [,4]
[1,]    1    6   11   16
[2,]    2    7   12   17
[3,]    3    8   13   18
[4,]    4    9   14   19
[5,]    5   10   15   20
\end{verbatim}

\begin{Shaded}
\begin{Highlighting}[]
\KeywordTok{class}\NormalTok{(m)}
\end{Highlighting}
\end{Shaded}

\begin{verbatim}
[1] "matrix"
\end{verbatim}

\begin{Shaded}
\begin{Highlighting}[]
\KeywordTok{typeof}\NormalTok{(m)}
\end{Highlighting}
\end{Shaded}

\begin{verbatim}
[1] "integer"
\end{verbatim}

\begin{Shaded}
\begin{Highlighting}[]
\CommentTok{# Dimensions of m}
\KeywordTok{dim}\NormalTok{(m) }
\end{Highlighting}
\end{Shaded}

\begin{verbatim}
[1] 5 4
\end{verbatim}

\hypertarget{list}{%
\subsubsection{List}\label{list}}

A `list' is a containter containing different objects. The contents of
list need not be of the same type or mode. A list can encompass a
mixture of data types such as vectors, matrices, data frames, other
lists or any other data structure.

\begin{Shaded}
\begin{Highlighting}[]
\NormalTok{w <-}\StringTok{ }\KeywordTok{list}\NormalTok{(a, m, d, }\KeywordTok{list}\NormalTok{(b, c))}
\KeywordTok{class}\NormalTok{(w)}
\end{Highlighting}
\end{Shaded}

\begin{verbatim}
[1] "list"
\end{verbatim}

\begin{Shaded}
\begin{Highlighting}[]
\KeywordTok{str}\NormalTok{(w)}
\end{Highlighting}
\end{Shaded}

\begin{verbatim}
List of 4
 $ : num [1:3] 1 2 3.3
 $ : int [1:5, 1:4] 1 2 3 4 5 6 7 8 9 10 ...
 $ : logi [1:6] TRUE TRUE TRUE FALSE TRUE FALSE
 $ :List of 2
  ..$ : int [1:3] 1 2 3
  ..$ : chr [1:3] "one" "two" "three"
\end{verbatim}

\hypertarget{dataframe}{%
\subsubsection{Data Frame}\label{dataframe}}

A `data frame' in \texttt{R} is a special kind of list with every
element having equal length. It is very important for handling tabular
data in \texttt{R}. It is a array like structure with rows and columns.
Each column needs to be of a single data type, however data type can
vary between columns.

\begin{Shaded}
\begin{Highlighting}[]
\NormalTok{L <-}\StringTok{ }\NormalTok{LETTERS[}\DecValTok{1}\OperatorTok{:}\DecValTok{4}\NormalTok{]}
\NormalTok{y <-}\StringTok{ }\DecValTok{1}\OperatorTok{:}\DecValTok{4}
\NormalTok{z <-}\StringTok{ }\KeywordTok{c}\NormalTok{(}\StringTok{"This"}\NormalTok{, }\StringTok{"is"}\NormalTok{, }\StringTok{"a"}\NormalTok{, }\StringTok{"data frame"}\NormalTok{)}
\NormalTok{df <-}\StringTok{ }\KeywordTok{data.frame}\NormalTok{(L, }\DataTypeTok{x =} \DecValTok{1}\NormalTok{, y, z)}
\NormalTok{df}
\end{Highlighting}
\end{Shaded}

\begin{verbatim}
  L x y          z
1 A 1 1       This
2 B 1 2         is
3 C 1 3          a
4 D 1 4 data frame
\end{verbatim}

\begin{Shaded}
\begin{Highlighting}[]
\KeywordTok{str}\NormalTok{(df)}
\end{Highlighting}
\end{Shaded}

\begin{verbatim}
'data.frame':   4 obs. of  4 variables:
 $ L: Factor w/ 4 levels "A","B","C","D": 1 2 3 4
 $ x: num  1 1 1 1
 $ y: int  1 2 3 4
 $ z: Factor w/ 4 levels "a","data frame",..: 4 3 1 2
\end{verbatim}

\begin{Shaded}
\begin{Highlighting}[]
\KeywordTok{attributes}\NormalTok{(df)}
\end{Highlighting}
\end{Shaded}

\begin{verbatim}
$names
[1] "L" "x" "y" "z"

$class
[1] "data.frame"

$row.names
[1] 1 2 3 4
\end{verbatim}

\begin{Shaded}
\begin{Highlighting}[]
\KeywordTok{rownames}\NormalTok{(df)}
\end{Highlighting}
\end{Shaded}

\begin{verbatim}
[1] "1" "2" "3" "4"
\end{verbatim}

\begin{Shaded}
\begin{Highlighting}[]
\KeywordTok{colnames}\NormalTok{(df)}
\end{Highlighting}
\end{Shaded}

\begin{verbatim}
[1] "L" "x" "y" "z"
\end{verbatim}

\hypertarget{functions}{%
\subsubsection{Functions}\label{functions}}

All of the work in \texttt{R} is done by functions. It is an object
defining a procedure which takes one or more objects as input (or
`arguments'), performs some action on them and finally gives a new
object as output (or `return'). \texttt{class()}, \texttt{mean()},
\texttt{getwd()}, \texttt{+}, etc. are all functions.

For example the function \texttt{mean()} takes a numeric vector as
argument and returns the mean as a numeric vector.

\begin{Shaded}
\begin{Highlighting}[]
\NormalTok{a <-}\StringTok{ }\KeywordTok{c}\NormalTok{(}\DecValTok{1}\NormalTok{, }\DecValTok{2}\NormalTok{, }\FloatTok{3.3}\NormalTok{)}
\KeywordTok{mean}\NormalTok{(a)}
\end{Highlighting}
\end{Shaded}

\begin{verbatim}
[1] 2.1
\end{verbatim}

The user can also create custom functions. For example the function
\texttt{foo} adds two numbers and gives the result.

\begin{Shaded}
\begin{Highlighting}[]
\NormalTok{foo <-}\StringTok{ }\ControlFlowTok{function}\NormalTok{(n1, n2) \{}
\NormalTok{  out <-}\StringTok{ }\NormalTok{n1 }\OperatorTok{+}\StringTok{ }\NormalTok{n2}
  \KeywordTok{return}\NormalTok{(out)}
\NormalTok{\}}
\KeywordTok{foo}\NormalTok{(}\DecValTok{2}\NormalTok{,}\DecValTok{3}\NormalTok{)}
\end{Highlighting}
\end{Shaded}

\begin{verbatim}
[1] 5
\end{verbatim}

\hypertarget{special-elements}{%
\subsection{Special Elements}\label{special-elements}}

In addition to numbers and text, there are some special elements which
can be included in different data objects.

\texttt{NA} (not available) indicates missing data.

\begin{Shaded}
\begin{Highlighting}[]
\NormalTok{x <-}\StringTok{ }\KeywordTok{c}\NormalTok{(}\FloatTok{2.5}\NormalTok{, }\OtherTok{NA}\NormalTok{, }\FloatTok{8.6}\NormalTok{)}
\NormalTok{y <-}\StringTok{ }\KeywordTok{c}\NormalTok{(}\OtherTok{TRUE}\NormalTok{, }\OtherTok{FALSE}\NormalTok{, }\OtherTok{NA}\NormalTok{)}
\NormalTok{z <-}\StringTok{ }\KeywordTok{c}\NormalTok{(}\StringTok{"k"}\NormalTok{, }\OtherTok{NA}\NormalTok{, }\StringTok{"m"}\NormalTok{, }\StringTok{"n"}\NormalTok{, }\StringTok{"o"}\NormalTok{)}
\KeywordTok{is.na}\NormalTok{(x)}
\end{Highlighting}
\end{Shaded}

\begin{verbatim}
[1] FALSE  TRUE FALSE
\end{verbatim}

\begin{Shaded}
\begin{Highlighting}[]
\KeywordTok{is.na}\NormalTok{(z)}
\end{Highlighting}
\end{Shaded}

\begin{verbatim}
[1] FALSE  TRUE FALSE FALSE FALSE
\end{verbatim}

\begin{Shaded}
\begin{Highlighting}[]
\KeywordTok{anyNA}\NormalTok{(x)}
\end{Highlighting}
\end{Shaded}

\begin{verbatim}
[1] TRUE
\end{verbatim}

\begin{Shaded}
\begin{Highlighting}[]
\NormalTok{a}
\end{Highlighting}
\end{Shaded}

\begin{verbatim}
[1] 1.0 2.0 3.3
\end{verbatim}

\begin{Shaded}
\begin{Highlighting}[]
\KeywordTok{is.na}\NormalTok{(a)}
\end{Highlighting}
\end{Shaded}

\begin{verbatim}
[1] FALSE FALSE FALSE
\end{verbatim}

\texttt{Inf} indicates infinity.

\begin{Shaded}
\begin{Highlighting}[]
\DecValTok{1}\OperatorTok{/}\DecValTok{0}
\end{Highlighting}
\end{Shaded}

\begin{verbatim}
[1] Inf
\end{verbatim}

\texttt{NaN} (Not a Number) indicates any undefined value.

\begin{Shaded}
\begin{Highlighting}[]
\DecValTok{0}\OperatorTok{/}\DecValTok{0}
\end{Highlighting}
\end{Shaded}

\begin{verbatim}
[1] NaN
\end{verbatim}

\hypertarget{indexing}{%
\subsection{Indexing}\label{indexing}}

The \texttt{{[}} function is used to extract elements of an object by
indexing (numeric or logical). Named elements in lists and data frames
can be extracted by using the \texttt{\$} operator.

Consider a vector \texttt{a}.

\begin{Shaded}
\begin{Highlighting}[]
\NormalTok{a <-}\StringTok{ }\KeywordTok{c}\NormalTok{(}\DecValTok{1}\NormalTok{, }\DecValTok{2}\NormalTok{, }\FloatTok{3.3}\NormalTok{, }\FloatTok{2.8}\NormalTok{, }\FloatTok{6.7}\NormalTok{)}
\CommentTok{# Numeric indexing}
\CommentTok{# Extract first element}
\NormalTok{a[}\DecValTok{1}\NormalTok{]}
\end{Highlighting}
\end{Shaded}

\begin{verbatim}
[1] 1
\end{verbatim}

\begin{Shaded}
\begin{Highlighting}[]
\CommentTok{# Extract elements 2:3}
\NormalTok{a[}\DecValTok{2}\OperatorTok{:}\DecValTok{3}\NormalTok{]}
\end{Highlighting}
\end{Shaded}

\begin{verbatim}
[1] 2.0 3.3
\end{verbatim}

\begin{Shaded}
\begin{Highlighting}[]
\CommentTok{# Logical indexing}
\NormalTok{a[a }\OperatorTok{>}\StringTok{ }\DecValTok{3}\NormalTok{]}
\end{Highlighting}
\end{Shaded}

\begin{verbatim}
[1] 3.3 6.7
\end{verbatim}

Consider a matrix \texttt{m}.

\begin{Shaded}
\begin{Highlighting}[]
\NormalTok{m <-}\StringTok{ }\KeywordTok{matrix}\NormalTok{(}\DecValTok{1}\OperatorTok{:}\DecValTok{9}\NormalTok{, }\DataTypeTok{nrow =} \DecValTok{3}\NormalTok{, }\DataTypeTok{ncol =} \DecValTok{3}\NormalTok{, }\DataTypeTok{byrow =} \OtherTok{TRUE}\NormalTok{)}
\KeywordTok{colnames}\NormalTok{(m) <-}\StringTok{ }\KeywordTok{c}\NormalTok{(}\StringTok{'a'}\NormalTok{, }\StringTok{'b'}\NormalTok{, }\StringTok{'c'}\NormalTok{)}
\NormalTok{m}
\end{Highlighting}
\end{Shaded}

\begin{verbatim}
     a b c
[1,] 1 2 3
[2,] 4 5 6
[3,] 7 8 9
\end{verbatim}

\begin{Shaded}
\begin{Highlighting}[]
\CommentTok{# Extract elements}
\NormalTok{m[,}\DecValTok{2}\NormalTok{] }\CommentTok{# 2nd column of matrix}
\end{Highlighting}
\end{Shaded}

\begin{verbatim}
[1] 2 5 8
\end{verbatim}

\begin{Shaded}
\begin{Highlighting}[]
\NormalTok{m[}\DecValTok{3}\NormalTok{,] }\CommentTok{# 3rd row of matrix}
\end{Highlighting}
\end{Shaded}

\begin{verbatim}
a b c 
7 8 9 
\end{verbatim}

\begin{Shaded}
\begin{Highlighting}[]
\NormalTok{m[}\DecValTok{2}\OperatorTok{:}\DecValTok{3}\NormalTok{, }\DecValTok{1}\OperatorTok{:}\DecValTok{3}\NormalTok{] }\CommentTok{# rows 2,3 of columns 1,2,3 }
\end{Highlighting}
\end{Shaded}

\begin{verbatim}
     a b c
[1,] 4 5 6
[2,] 7 8 9
\end{verbatim}

\begin{Shaded}
\begin{Highlighting}[]
\NormalTok{m[}\DecValTok{2}\NormalTok{,}\DecValTok{2}\NormalTok{] }\CommentTok{# Element in 2nd column of 2nd row}
\end{Highlighting}
\end{Shaded}

\begin{verbatim}
b 
5 
\end{verbatim}

\begin{Shaded}
\begin{Highlighting}[]
\NormalTok{m[, }\StringTok{'b'}\NormalTok{] }\CommentTok{# Column 'b'}
\end{Highlighting}
\end{Shaded}

\begin{verbatim}
[1] 2 5 8
\end{verbatim}

\begin{Shaded}
\begin{Highlighting}[]
\NormalTok{m[, }\KeywordTok{c}\NormalTok{(}\StringTok{'a'}\NormalTok{, }\StringTok{'c'}\NormalTok{)] }\CommentTok{# Column 'a' and 'c'}
\end{Highlighting}
\end{Shaded}

\begin{verbatim}
     a c
[1,] 1 3
[2,] 4 6
[3,] 7 9
\end{verbatim}

Consider a list \texttt{w}.

\begin{Shaded}
\begin{Highlighting}[]
\NormalTok{w <-}\StringTok{ }\KeywordTok{list}\NormalTok{(}\DataTypeTok{vec =}\NormalTok{ a, }\DataTypeTok{mat =}\NormalTok{ m, }\DataTypeTok{data =}\NormalTok{ df, }\DataTypeTok{alist =} \KeywordTok{list}\NormalTok{(b, c))}

\CommentTok{# Indexing by number}
\NormalTok{w[}\DecValTok{2}\NormalTok{] }\CommentTok{# As list structure}
\end{Highlighting}
\end{Shaded}

\begin{verbatim}
$mat
     a b c
[1,] 1 2 3
[2,] 4 5 6
[3,] 7 8 9
\end{verbatim}

\begin{Shaded}
\begin{Highlighting}[]
\NormalTok{w[[}\DecValTok{2}\NormalTok{]] }\CommentTok{# Without list structure}
\end{Highlighting}
\end{Shaded}

\begin{verbatim}
     a b c
[1,] 1 2 3
[2,] 4 5 6
[3,] 7 8 9
\end{verbatim}

\begin{Shaded}
\begin{Highlighting}[]
\CommentTok{# Indexing by name}
\NormalTok{w}\OperatorTok{$}\NormalTok{vec}
\end{Highlighting}
\end{Shaded}

\begin{verbatim}
[1] 1.0 2.0 3.3 2.8 6.7
\end{verbatim}

\begin{Shaded}
\begin{Highlighting}[]
\NormalTok{w}\OperatorTok{$}\NormalTok{data}
\end{Highlighting}
\end{Shaded}

\begin{verbatim}
  L x y          z
1 A 1 1       This
2 B 1 2         is
3 C 1 3          a
4 D 1 4 data frame
\end{verbatim}

Consider a data frame \texttt{df}.

\begin{Shaded}
\begin{Highlighting}[]
\NormalTok{df}
\end{Highlighting}
\end{Shaded}

\begin{verbatim}
  L x y          z
1 A 1 1       This
2 B 1 2         is
3 C 1 3          a
4 D 1 4 data frame
\end{verbatim}

\begin{Shaded}
\begin{Highlighting}[]
\CommentTok{# Indexing by number}
\NormalTok{df[,}\DecValTok{2}\NormalTok{] }\CommentTok{# 2nd column of data frame}
\end{Highlighting}
\end{Shaded}

\begin{verbatim}
[1] 1 1 1 1
\end{verbatim}

\begin{Shaded}
\begin{Highlighting}[]
\NormalTok{df[}\DecValTok{2}\NormalTok{] }\CommentTok{# 2nd column of data frame}
\end{Highlighting}
\end{Shaded}

\begin{verbatim}
  x
1 1
2 1
3 1
4 1
\end{verbatim}

\begin{Shaded}
\begin{Highlighting}[]
\NormalTok{df[}\DecValTok{3}\NormalTok{,] }\CommentTok{# 3rd row of data frame}
\end{Highlighting}
\end{Shaded}

\begin{verbatim}
  L x y z
3 C 1 3 a
\end{verbatim}

\begin{Shaded}
\begin{Highlighting}[]
\NormalTok{df[}\DecValTok{2}\OperatorTok{:}\DecValTok{3}\NormalTok{, }\DecValTok{1}\OperatorTok{:}\DecValTok{3}\NormalTok{] }\CommentTok{# rows 2,3 of columns 1,2,3 }
\end{Highlighting}
\end{Shaded}

\begin{verbatim}
  L x y
2 B 1 2
3 C 1 3
\end{verbatim}

\begin{Shaded}
\begin{Highlighting}[]
\NormalTok{df[}\DecValTok{2}\NormalTok{,}\DecValTok{2}\NormalTok{] }\CommentTok{# Element in 2nd column of 2nd row}
\end{Highlighting}
\end{Shaded}

\begin{verbatim}
[1] 1
\end{verbatim}

\begin{Shaded}
\begin{Highlighting}[]
\CommentTok{# Indexing by name}
\NormalTok{df}\OperatorTok{$}\NormalTok{L}
\end{Highlighting}
\end{Shaded}

\begin{verbatim}
[1] A B C D
Levels: A B C D
\end{verbatim}

\begin{Shaded}
\begin{Highlighting}[]
\NormalTok{df}\OperatorTok{$}\NormalTok{z}
\end{Highlighting}
\end{Shaded}

\begin{verbatim}
[1] This       is         a          data frame
Levels: a data frame is This
\end{verbatim}

\hypertarget{help-documentation}{%
\subsection{Help Documentation}\label{help-documentation}}

The help documentation regarding any function can be viewed using the
\texttt{?} or \texttt{help()} function. The help documentation shows the
default usage of the function including, the arguments that are taken by
the function and the type of output object returned (`Value').

\begin{Shaded}
\begin{Highlighting}[]
\NormalTok{?ls}
\KeywordTok{help}\NormalTok{(ls)}

\NormalTok{?mean}

\NormalTok{?setwd}
\end{Highlighting}
\end{Shaded}

\hypertarget{pack}{%
\subsection{Packages}\label{pack}}

Packages in \texttt{R} are collections of \texttt{R} functions, data,
and compiled code in a well-defined format. They are add-ons which
extend the functionality of \texttt{R} and at present, there are
\href{https://cran.r-project.org/web/packages/available_packages_by_name.html}{13646}
packages available for deployment and use at the official repository,
the Comprehensive R Archive Network (CRAN).

Valid packages from CRAN can be installed by using the
\texttt{install.packages()} command.

\begin{Shaded}
\begin{Highlighting}[]
\CommentTok{# Install the package 'readxl' for importing data from excel}
\KeywordTok{install.packages}\NormalTok{(readxl)}
\end{Highlighting}
\end{Shaded}

Installed packages can be loaded using the function \texttt{library()}.

\begin{Shaded}
\begin{Highlighting}[]
\CommentTok{# Install the package 'readxl' for importing data from excel}
\KeywordTok{library}\NormalTok{(readxl)}
\end{Highlighting}
\end{Shaded}

\hypertarget{impexp}{%
\subsection{Importing and Exporting Tabular Data}\label{impexp}}

Tabular data from a spreadsheet can be imported into \texttt{R} in
different ways. Consider some data such as in Table 1. Copy this data in
to a spreadsheet editor such as MS Excel and save it as
\texttt{augdata.csv}, a comma-separated-value file and
\texttt{augdata.xlsx}, an Excel file in the working directory
(\texttt{getwd()}).

\begin{center}
                \textbf{Table 1}: Example data from an experiment in augmented RCBD design.
                \end{center}

\begin{longtable}[]{@{}lrrr@{}}
\toprule
blk & trt & y1 & y2\tabularnewline
\midrule
\endhead
I & 1 & 92 & 258\tabularnewline
I & 2 & 79 & 224\tabularnewline
I & 3 & 87 & 238\tabularnewline
I & 4 & 81 & 278\tabularnewline
I & 7 & 96 & 347\tabularnewline
I & 11 & 89 & 300\tabularnewline
I & 12 & 82 & 289\tabularnewline
II & 1 & 79 & 260\tabularnewline
II & 2 & 81 & 220\tabularnewline
II & 3 & 81 & 237\tabularnewline
II & 4 & 91 & 227\tabularnewline
II & 5 & 79 & 281\tabularnewline
II & 9 & 78 & 311\tabularnewline
III & 1 & 83 & 250\tabularnewline
III & 2 & 77 & 240\tabularnewline
III & 3 & 78 & 268\tabularnewline
III & 4 & 78 & 287\tabularnewline
III & 8 & 70 & 226\tabularnewline
III & 6 & 75 & 395\tabularnewline
III & 10 & 74 & 450\tabularnewline
\bottomrule
\end{longtable}

The \texttt{augdata.csv} file can be imported into \texttt{R} using the
\texttt{read.csv()} function or the
\href{https://readr.tidyverse.org/reference/read_delim.html}{\texttt{read\_csv()}}
function in the \texttt{readr} package.

\begin{Shaded}
\begin{Highlighting}[]
\NormalTok{data <-}\StringTok{ }\KeywordTok{read.csv}\NormalTok{(}\DataTypeTok{file =} \StringTok{"augdata.csv"}\NormalTok{)}
\KeywordTok{str}\NormalTok{(data)}
\end{Highlighting}
\end{Shaded}

\begin{verbatim}
'data.frame':   20 obs. of  4 variables:
 $ blk: Factor w/ 3 levels "I","II","III": 1 1 1 1 1 1 1 2 2 2 ...
 $ trt: num  1 2 3 4 7 11 12 1 2 3 ...
 $ y1 : num  92 79 87 81 96 89 82 79 81 81 ...
 $ y2 : num  258 224 238 278 347 300 289 260 220 237 ...
\end{verbatim}

The argument \texttt{stringsAsFactors\ =\ FALSE} reads the text columns
as of type \texttt{character} instead of the default \texttt{factor}.

\begin{Shaded}
\begin{Highlighting}[]
\NormalTok{data <-}\StringTok{ }\KeywordTok{read.csv}\NormalTok{(}\DataTypeTok{file =} \StringTok{"augdata.csv"}\NormalTok{, }\DataTypeTok{stringsAsFactors =} \OtherTok{FALSE}\NormalTok{)}
\KeywordTok{str}\NormalTok{(data)}
\end{Highlighting}
\end{Shaded}

\begin{verbatim}
'data.frame':   20 obs. of  4 variables:
 $ blk: chr  "I" "I" "I" "I" ...
 $ trt: num  1 2 3 4 7 11 12 1 2 3 ...
 $ y1 : num  92 79 87 81 96 89 82 79 81 81 ...
 $ y2 : num  258 224 238 278 347 300 289 260 220 237 ...
\end{verbatim}

The \texttt{augdata.xlsx} file can be imported into \texttt{R} using the
\href{https://readxl.tidyverse.org/reference/read_excel.html}{\texttt{read\_excel()}}
function in the \texttt{readxl} package.

\begin{Shaded}
\begin{Highlighting}[]
\KeywordTok{library}\NormalTok{(readxl)}
\NormalTok{data <-}\StringTok{ }\KeywordTok{read_excel}\NormalTok{(}\DataTypeTok{path =} \StringTok{"augdata.xlsx"}\NormalTok{)}
\end{Highlighting}
\end{Shaded}

\begin{verbatim}
'data.frame':   20 obs. of  4 variables:
 $ blk: chr  "I" "I" "I" "I" ...
 $ trt: num  1 2 3 4 7 11 12 1 2 3 ...
 $ y1 : num  92 79 87 81 96 89 82 79 81 81 ...
 $ y2 : num  258 224 238 278 347 300 289 260 220 237 ...
\end{verbatim}

The tabular data can be exported from \texttt{R} to a \texttt{.csv}
(comma-separated-value) file by the
\href{https://www.rdocumentation.org/packages/utils/versions/3.5.1/topics/write.table}{\texttt{write.csv()}}
function.

\begin{Shaded}
\begin{Highlighting}[]
\KeywordTok{write.csv}\NormalTok{(}\DataTypeTok{x =}\NormalTok{ data, }\DataTypeTok{file =} \StringTok{"augdata.csv"}\NormalTok{)}
\end{Highlighting}
\end{Shaded}

\hypertarget{additional-resources}{%
\subsection{Additional Resources}\label{additional-resources}}

To learn more about \texttt{R}, there are upteen number of online
tutorials as well as free courses availble. Queries about various
aspects can be put to the active and vibrant `R community online.

\begin{itemize}
\tightlist
\item
  Online tutorials

  \begin{itemize}
  \tightlist
  \item
    \url{http://www.cran.r-project.org/other-docs.html}
  \item
    \url{https://bookdown.org/ndphillips/YaRrr/}
  \end{itemize}
\item
  Free online courses

  \begin{itemize}
  \tightlist
  \item
    \url{http://tryr.codeschool.com/}
  \item
    \url{https://www.datacamp.com/courses/free-introduction-to-r}
  \end{itemize}
\item
  \texttt{R} community support

  \begin{itemize}
  \tightlist
  \item
    \url{http://stackoverflow.com/}
  \item
    \texttt{R} help mailing lists :
    \url{http://www.r-project.org/mail.html}
  \end{itemize}
\end{itemize}

\hypertarget{install}{%
\section{\texorpdfstring{Installation of
\texttt{augmentedRCBD}}{Installation of augmentedRCBD}}\label{install}}

The package \texttt{augmentedRCBD} can be installed using the following
functions:

\begin{Shaded}
\begin{Highlighting}[]
\CommentTok{# Install from CRAN}
\KeywordTok{install.packages}\NormalTok{(}\StringTok{'augmentedRCBD'}\NormalTok{, }\DataTypeTok{dependencies=}\OtherTok{TRUE}\NormalTok{)}

\CommentTok{# Install development version from Github}

\KeywordTok{library}\NormalTok{(devtools)}
\KeywordTok{install_github}\NormalTok{(}\StringTok{"aravind-j/augmentedRCBD"}\NormalTok{)}
\end{Highlighting}
\end{Shaded}

The stable release is hosted in
\href{https://CRAN.R-project.org/package=augmentedRCBD}{CRAN}
\protect\hyperlink{pack}{(see section 4.7)}, while the under-development
version is hosted as a
\href{https://github.com/aravind-j/augmentedRCBD}{Github} repository. To
install from github, you need to use the
\href{https://devtools.r-lib.org/reference/reexports.html}{\texttt{install\_github()}}
function from \href{https://devtools.r-lib.org/}{`devtools} package.

\hypertarget{data-format}{%
\section{Data Format}\label{data-format}}

Certain details need to be considered for arranging experimental data
for analysis using the \texttt{augmentedRCBD} package.

The data should be in long/vertical form, where each row has the data
from one genotype per block. For example, consider the following data
(Table 2) recorded for a trait from an experiment laid out in an
augmented block design with 3 blocks and 12 genotypes(or treatment) with
6 to 7 genotypes/block. 8 genotypes (Test, G 5 to G 12) are not
replicated, while 4 genotypes (Check, G 1 to G 4) are replicated.

\begin{center}
                \textbf{Table 2}: Data from an experiment in augmented RCBD design.
                \end{center}

\begin{longtable}[]{@{}llllllll@{}}
\toprule
\endhead
\textbf{Block I} & G12 & \textbf{G4} & G11 & \textbf{G2} & \textbf{G1} &
G7 & \textbf{G3}\tabularnewline
& 82 & 81 & 89 & 79 & 92 & 96 & 87\tabularnewline
\textbf{Block II} & G5 & G9 & -- & \textbf{G3} & \textbf{G1} &
\textbf{G2} & \textbf{G4}\tabularnewline
& 79 & 78 & -- & 81 & 79 & 81 & 91\tabularnewline
\textbf{Block III} & \textbf{G4} & \textbf{G2} & \textbf{G1} & G6 & G10
& \textbf{G3} & G8\tabularnewline
& 78 & 77 & 83 & 75 & 74 & 78 & 70\tabularnewline
\bottomrule
\end{longtable}

This data needs to be arranged with columns showing block, genotype (or
treatment) and the data of the trait for each genotype per block (Table
3).

\begin{center}
                \textbf{Table 3}: Data from an experiment in augmented RCBD design arranged in long-form.
                \end{center}

\begin{longtable}[]{@{}llr@{}}
\toprule
Block & Treatment & Trait\tabularnewline
\midrule
\endhead
Block I & G 1 & 92\tabularnewline
Block I & G 2 & 79\tabularnewline
Block I & G 3 & 87\tabularnewline
Block I & G 4 & 81\tabularnewline
Block I & G 7 & 96\tabularnewline
Block I & G 11 & 89\tabularnewline
Block I & G 12 & 82\tabularnewline
Block II & G 1 & 79\tabularnewline
Block II & G 2 & 81\tabularnewline
Block II & G 3 & 81\tabularnewline
Block II & G 4 & 91\tabularnewline
Block II & G 5 & 79\tabularnewline
Block II & G 9 & 78\tabularnewline
Block III & G 1 & 83\tabularnewline
Block III & G 2 & 77\tabularnewline
Block III & G 3 & 78\tabularnewline
Block III & G 4 & 78\tabularnewline
Block III & G 8 & 70\tabularnewline
Block III & G 6 & 75\tabularnewline
Block III & G 10 & 74\tabularnewline
\bottomrule
\end{longtable}

The data for block and genotype (or treatment) can also be depicted as
numbers (Table 4).

\begin{center}
                \textbf{Table 4}: Data from an experiment in augmented RCBD design arranged in long-form (Block and Treatment as numbers).
                \end{center}

\begin{longtable}[]{@{}rrr@{}}
\toprule
Block & Treatment & Trait\tabularnewline
\midrule
\endhead
1 & 1 & 92\tabularnewline
1 & 2 & 79\tabularnewline
1 & 3 & 87\tabularnewline
1 & 4 & 81\tabularnewline
1 & 7 & 96\tabularnewline
1 & 11 & 89\tabularnewline
1 & 12 & 82\tabularnewline
2 & 1 & 79\tabularnewline
2 & 2 & 81\tabularnewline
2 & 3 & 81\tabularnewline
2 & 4 & 91\tabularnewline
2 & 5 & 79\tabularnewline
2 & 9 & 78\tabularnewline
3 & 1 & 83\tabularnewline
3 & 2 & 77\tabularnewline
3 & 3 & 78\tabularnewline
3 & 4 & 78\tabularnewline
3 & 8 & 70\tabularnewline
3 & 6 & 75\tabularnewline
3 & 10 & 74\tabularnewline
\bottomrule
\end{longtable}

Multiple traits can be added as additional columns (Table 5).

\begin{center}
                \textbf{Table 5}: Data from an experiment in augmented RCBD design arranged in long-form (Multiple traits).
                \end{center}

\begin{longtable}[]{@{}llrr@{}}
\toprule
Block & Treatment & Trait1 & Trait2\tabularnewline
\midrule
\endhead
Block I & G 1 & 92 & 258\tabularnewline
Block I & G 2 & 79 & 224\tabularnewline
Block I & G 3 & 87 & 238\tabularnewline
Block I & G 4 & 81 & 278\tabularnewline
Block I & G 7 & 96 & 347\tabularnewline
Block I & G 11 & 89 & 300\tabularnewline
Block I & G 12 & 82 & 289\tabularnewline
Block II & G 1 & 79 & 260\tabularnewline
Block II & G 2 & 81 & 220\tabularnewline
Block II & G 3 & 81 & 237\tabularnewline
Block II & G 4 & 91 & 227\tabularnewline
Block II & G 5 & 79 & 281\tabularnewline
Block II & G 9 & 78 & 311\tabularnewline
Block III & G 1 & 83 & 250\tabularnewline
Block III & G 2 & 77 & 240\tabularnewline
Block III & G 3 & 78 & 268\tabularnewline
Block III & G 4 & 78 & 287\tabularnewline
Block III & G 8 & 70 & 226\tabularnewline
Block III & G 6 & 75 & 395\tabularnewline
Block III & G 10 & 74 & 450\tabularnewline
\bottomrule
\end{longtable}

Data should preferably be balanced i.e.~all the check genotypes should
be present in all the blocks. If not, a warning is issued. The number of
test genotypes can vary within a block. There should not be any missing
values. Rows of genotypes with missing values for one or more traits
should be removed.

Such a tabular data should be imported \protect\hyperlink{impexp}{see
section 7.8} into \texttt{R} as a data frame object
\protect\hyperlink{dataframe}{see section 4.3.5}. The columns with the
block and treatment categorical data should of the type factor
\protect\hyperlink{factor}{see section 4.3.2}, while the column(s) with
the trait data should be of the type integer or numeric
\protect\hyperlink{vector}{see section 4.3.1}.

Consider the example data in Table 1 in
\protect\hyperlink{impexp}{section 4.8}. It can be imported as
instructed.

\hypertarget{data-analysis-for-a-single-trait}{%
\section{Data Analysis for a Single
Trait}\label{data-analysis-for-a-single-trait}}

Check genotypes are inferred by default on the basis of number of
replications. However, if some test genotypes are also replicated, they
may also be falsely detected as checks.

\includegraphics{C:/Users/lenovo/Dropbox/RPACKA~1/AUGMEN~1/docs/articles/DATA_A~1/figure-latex/unnamed-chunk-56-1.pdf}

\textbf{Fig. 4}. Workflow for analysis of single traits with
\texttt{augmentedRCBD}.

\hypertarget{data-analysis-for-a-multiple-traits}{%
\section{Data Analysis for a Multiple
Traits}\label{data-analysis-for-a-multiple-traits}}

\includegraphics{C:/Users/lenovo/Dropbox/RPACKA~1/AUGMEN~1/docs/articles/DATA_A~1/figure-latex/unnamed-chunk-57-1.pdf}

\textbf{Fig. 5}. Workflow for analysis of multiple traits with
\texttt{augmentedRCBD}.

\clearpage

Then the package can be loaded using the function

\begin{Shaded}
\begin{Highlighting}[]
\KeywordTok{library}\NormalTok{(augmentedRCBD)}
\end{Highlighting}
\end{Shaded}

\begin{verbatim}
Registered S3 methods overwritten by 'ggplot2':
  method         from 
  [.quosures     rlang
  c.quosures     rlang
  print.quosures rlang
\end{verbatim}

\begin{verbatim}
Registered S3 method overwritten by 'dplyr':
  method               from  
  as.data.frame.tbl_df tibble
\end{verbatim}

\begin{verbatim}

--------------------------------------------------------------------------------
Welcome to augmentedRCBD version 0.1.0


# To know whats new in this version type:
  news(package='augmentedRCBD')
  for the NEWS file.

# To cite the methods in the package type:
  citation(package='augmentedRCBD')

# To suppress this message use:
  suppressPackageStartupMessages(library(augmentedRCBD))
--------------------------------------------------------------------------------
\end{verbatim}

\hypertarget{citing-augmentedrcbd}{%
\section{\texorpdfstring{Citing
\texttt{augmentedRCBD}}{Citing augmentedRCBD}}\label{citing-augmentedrcbd}}

\begin{verbatim}

To cite the R package 'augmentedRCBD' in publications use:

  J. Aravind, Mukesh Sankar S., Dhammaprakash P. Wankhede, and
  Vikender Kaur (2018).  augmentedRCBD: Analysis of Augmented
  Randomised Complete Block Designs. R package version 0.1.0,
  https://aravind-j.github.io/augmentedRCBD/.

A BibTeX entry for LaTeX users is

  @Manual{,
    title = {augmentedRCBD: Analysis of Augmented Randomised Complete Block Designs},
    author = {{Aravind J} and {Mukesh Sankar S} and {Dhammaprakash Pandhari Wankhede} and {Vikender Kaur}},
    year = {2018},
    note = {R package version 0.1.0},
    note = {https://aravind-j.github.io/augmentedRCBD/},
  }

This free and open-source software implements academic research by
the authors and co-workers. If you use it, please support the
project by citing the package.
\end{verbatim}

\hypertarget{session-info}{%
\section{Session Info}\label{session-info}}

\begin{Shaded}
\begin{Highlighting}[]
\KeywordTok{sessionInfo}\NormalTok{()}
\end{Highlighting}
\end{Shaded}

\begin{verbatim}
R Under development (unstable) (2018-10-21 r75476)
Platform: i386-w64-mingw32/i386 (32-bit)
Running under: Windows 7 (build 7601) Service Pack 1

Matrix products: default

locale:
[1] LC_COLLATE=English_India.1252  LC_CTYPE=English_India.1252   
[3] LC_MONETARY=English_India.1252 LC_NUMERIC=C                  
[5] LC_TIME=English_India.1252    

attached base packages:
[1] stats     graphics  grDevices utils     datasets  methods   base     

other attached packages:
[1] augmentedRCBD_0.1.0 diagram_1.6.4       shape_1.4.4        

loaded via a namespace (and not attached):
 [1] bibtex_0.4.2       reshape2_1.4.3     magrittr_1.5      
 [4] TH.data_1.0-9      estimability_1.3   gtable_0.2.0      
 [7] rmarkdown_1.10     base64enc_0.1-3    xml2_1.2.0        
[10] pillar_1.3.0       htmltools_0.3.6    stringr_1.3.1     
[13] splines_3.6.0      dplyr_0.7.7        moments_0.14      
[16] lattice_0.20-35    survival_2.42-6    tidyselect_0.2.5  
[19] plyr_1.8.4         sandwich_2.5-0     emmeans_1.2.4     
[22] zoo_1.8-4          uuid_0.1-2         knitr_1.20        
[25] gbRd_0.4-11        pkgconfig_2.0.2    Matrix_1.2-14     
[28] R6_2.3.0           digest_0.6.18      colorspace_1.3-2  
[31] bindrcpp_0.2.2     rprojroot_1.3-2    stringi_1.2.4     
[34] yaml_2.2.0         lazyeval_0.2.1     codetools_0.2-15  
[37] evaluate_0.12      officer_0.3.2      tibble_1.4.2      
[40] compiler_3.6.0     gdtools_0.1.7      bindr_0.1.1       
[43] multcompView_0.1-7 backports_1.1.2    xtable_1.8-3      
[46] Rdpack_0.10-3      munsell_0.5.0      Rcpp_1.0.0        
[49] highr_0.7          coda_0.19-2        MASS_7.3-51       
[52] assertthat_0.2.0   ggplot2_3.0.0      tools_3.6.0       
[55] mvtnorm_1.0-8      zip_1.0.0          scales_1.0.0      
[58] crayon_1.3.4       glue_1.3.0         purrr_0.2.5       
[61] flextable_0.4.5    rlang_0.3.0.1      multcomp_1.4-8    
[64] grid_3.6.0        
\end{verbatim}

\hypertarget{references}{%
\section*{References}\label{references}}
\addcontentsline{toc}{section}{References}

\hypertarget{refs}{}
\leavevmode\hypertarget{ref-federer_augmented_1956}{}%
Federer, W. T. (1956). Augmented (or hoonuiaku) designs. \emph{The
Hawaiian Planters' Record} LV(2), 191--208.

\leavevmode\hypertarget{ref-federer_augmented_1961}{}%
Federer, W. T. (1961). Augmented designs with one-way elimination of
heterogeneity. \emph{Biometrics} 17, 447--473.
doi:\href{https://doi.org/10.2307/2527837}{10.2307/2527837}.

\leavevmode\hypertarget{ref-tippmann_programming_2015}{}%
Tippmann, S. (2015). Programming tools: Adventures with R. \emph{Nature
News} 517, 109.
doi:\href{https://doi.org/10.1038/517109a}{10.1038/517109a}.
  </div>

  <div class="col-md-3 hidden-xs hidden-sm" id="sidebar">
      </div>

</div>


      <footer>
      <div class="copyright">
  <p>Developed by J. Aravind, S. Mukesh Sankar, Dhammaprakash Pandhari Wankhede, Vikender Kaur.</p>
</div>

<div class="pkgdown">
  <p>Site built with <a href="http://pkgdown.r-lib.org/">pkgdown</a>.</p>
</div>

      </footer>
   </div>

  
<script type="text/javascript" src="https://cdn.jsdelivr.net/npm/docsearch.js@2/dist/cdn/docsearch.min.js"></script>
<script>
  docsearch({
    
    
    apiKey: 'bdafb7ea4490644e670f470ee3ef537f',
    indexName: 'augmentedrcbd',
    inputSelector: 'input#search-input.form-control',
    transformData: function(hits) {
      return hits.map(function (hit) {
        hit.url = updateHitURL(hit);
        return hit;
      });
    }
  });
</script>


  </body>
</html>

